\section{Related Work}
\label{sec:related_work}

This section reviews the literature most relevant to the themes developed in this
monograph. The discussion begins with formal methods for task and motion
synthesis, then turns to planning under incomplete models and infeasible
specifications, multi-robot coordination under temporal tasks, and execution-time
adaptation under interaction constraints. The final part focuses on contingent
collaboration and human supervision, which have become increasingly important in
large-scale and long-horizon robotic systems.

\subsection{Formal methods for task and motion synthesis}

Model checking has provided a principled foundation for synthesizing control and
planning policies from high-level specifications. Foundational work on symbolic
abstraction and discrete verification established the algorithmic basis for
reasoning about dynamical systems through finite models
\cite{alur2000discrete,girard2007approximately,baier2008principles}. Building on
this foundation, temporal-logic-based control synthesis has been developed for
several system classes, including discrete-time linear systems, piecewise affine
systems, and stochastic models
\cite{tabuada2006linear,aydin2013temporal,yordanov2010formal,
  kloetzer2008fully,yordanov2012temporal,ding2011mdp}. Metric-temporal and
robustness-based variants further extended this line of work beyond Boolean task
satisfaction by quantifying the degree of conformity of a trajectory to the given
specification \cite{fainekos2009robustness,abbas2013probabilistic}.

Within robotics, these developments led to a widely used architecture in which a
continuous robot model is abstracted into a finite transition system over a
partitioned workspace, after which a discrete plan is synthesized from a temporal
formula and refined into executable motion. Early robotic formulations of this
type were developed for autonomous motion planning under linear temporal logic
(LTL) tasks \cite{fainekos2009temporal}. Related efforts also considered user
interfaces for specifying temporal tasks and hierarchical structures for
manipulation and actuation \cite{srinivas2013graphical,he2015towards}. When the
mission includes region-dependent actions in addition to motion, integrated task
and motion planning becomes necessary. Existing formulations often rely either on
actions that are permanently attached to locations or on action propositions that
can be toggled independently of state, which is restrictive when actions depend
on local conditions, resource availability, or discrete choices among feasible
alternatives \cite{smith2011optimal,kress2009temporal}. This limitation motivates
the motion--action coupling developed later in Chapter~\ref{chapter:action},
which is closely related to the framework in \cite{guo2013motion}.

\subsection{Planning under incomplete models and infeasible specifications}

A major limitation of classical temporal-logic synthesis is the assumption that
the workspace and its relevant semantic properties are known accurately before
planning begins. Under this assumption, plans are typically produced offline and
executed against a fixed model, which reduces responsiveness to unexpected
changes. Several approaches relax this requirement by modeling environmental
uncertainty explicitly. Markov decision process formulations and robust planning
methods address uncertain observations and transitions probabilistically
\cite{ding2011mdp,wolff2012robust}, while GR(1)-style game formulations and
receding-horizon schemes embed adversarial or reactive environment behavior into
the synthesis loop \cite{kress2009temporal,wongpiromsarn2010receding}.
Perception-aware temporal planning has pushed this direction further by coupling
formal synthesis with uncertain semantic maps \cite{kantaros2022perception}.

A complementary line of work addresses model mismatch at execution time through
plan repair and revision. Local backtracking and transition patching have been
used to repair invalid discrete plans online, although such methods may not fully
capture changes in the truth values of workspace propositions
\cite{livingston2012backtracking}. Iterative refinement of the workspace model
from sensory updates has also been studied for co-safe temporal tasks
\cite{maly2013iterative}. The plan-revision perspective developed in
Chapter~\ref{chapter:partially}, based on \cite{guo2013revising}, belongs to this
broader class of approaches, but places particular emphasis on preserving formal
correctness while incorporating newly acquired information during execution.

A related challenge arises when the nominal task itself is unrealizable or
infeasible. In such cases, simply reporting failure offers limited guidance to
the user or the supervisory layer. Relaxation-based approaches search for the
closest realizable specification automaton or otherwise identify minimal changes
that restore feasibility \cite{fainekos2011revising,kim2012approximate}. Other
work analyzes the source of unrealizability by isolating conflicting environment
or system assumptions \cite{raman2011analyzing}. In parallel, least-violating or
maximally satisfying planning methods quantify the extent to which a task can be
satisfied when full realization is impossible
\cite{tumova2013least,tumova2013minimum,lahijanian2015time,
  tumova2014maximally}. The treatment developed in Chapter~\ref{chapter:partially}
is closest to \cite{guo2012infeasible,guo2015multi}, where hard and soft
constraints are distinguished explicitly and task violation is assessed not only
by occurrence but also by severity.

\subsection{Temporal-logic coordination in multi-robot systems}

The single-robot planning framework described above has been extended in several
directions to cooperative multi-robot systems. Early work on formal
multi-robot coordination mainly adopted a top-down viewpoint: a global temporal
task is specified for the team, decomposed into synchronized or partially
synchronized local behaviors, and then executed through coordinated motion
\cite{kloetzer2010automatic,ding2011automatic}. Robustness with respect to travel
time uncertainty and optimal path construction under global temporal constraints
have also been studied \cite{ulusoy2012robust,ulusoy2013optimality}. Mixed-integer
programming formulations similarly provide a direct encoding of routing and
synchronization constraints for temporal missions \cite{karaman2011sampling}.
Although these methods offer strong guarantees, they usually rely on centralized
synthesis and can become computationally demanding as the number of robots,
synchronization points, and temporal couplings increases.

More recent work has broadened the scope of temporal-logic coordination by
integrating task allocation, hierarchy, and heterogeneity into the synthesis
process. Surveys on multi-robot systems, cooperative planning, and distributed
constraint optimization highlight the importance of scalable assignment and
coalition mechanisms in heterogeneous robotic teams
\cite{arai2002advances,torreno2017cooperative,fioretto2018distributed}. In this
spirit, market-based and learning-enabled allocation methods have been proposed
for uncertain task rewards and dynamic coalition formation
\cite{zhao2021market,dai2024dynamic}. For temporal-logic missions in particular,
simultaneous task allocation and planning, sampling-based optimal synthesis,
abstraction-free methods, counting logics, and scalable coordination frameworks
have considerably expanded the size and complexity of solvable problems
\cite{schillinger2018simultaneous,kantaros2020stylus,luo2021abstraction,
  luo2022temporal,sahin2019multirobot,leahy2021scalable,cardona2024planning}.
Hierarchical decompositions and large-scale continual coordination have further
improved tractability for structured missions
\cite{luo2024decomposition,luo2025hulk,luo2025simultaneous}.

Alongside these top-down and team-level formulations, another thread considers
locally assigned specifications and loosely coupled collaboration. This bottom-up
view is particularly relevant when robots have heterogeneous capabilities,
asynchronous execution, and only intermittent dependence on one another. Local
plan reconfiguration, distributed conflict resolution, and coordination of task
and motion under local specifications have been studied in
\cite{guo2014distributed,guo2015multi,guo2015bottom,guo2016task}. Related
frameworks address decentralized feasibility verification and receding-horizon
coordination for dependent local temporal tasks
\cite{filippidis2012decentralized,tumova2014receding}. The material developed in
Chapters~\ref{chapter:partially}--\ref{chapter:continuous} is positioned most
naturally within this locally specified and loosely coupled branch of the
literature.

\subsection{Relative-motion and interaction constraints}

Formal task satisfaction alone is often insufficient for safe multi-robot
operation, because robots must also respect interaction constraints induced by
shared physical space and local communication structure. Typical examples include
collision avoidance, network connectivity, and relative-motion or relative-speed
constraints \cite{dimarogonas2006feedback,ji2007distributed,zavlanos2011graph,
  guo2014controlling}. These requirements are frequently imposed as standing
assumptions, even though they play a central role in the integrity and safety of
the team. Dynamic leader--follower coordination and decentralized control laws
have therefore been developed to enforce relative-distance constraints while
supporting multiple active task leaders and heterogeneous local goals
\cite{guo2014cooperative,guo2015hybrid,guo2015communication,kloetzer2011multi}.
This perspective is closely connected to the continuous-control developments in
Chapter~\ref{chapter:continuous}.

Embedded graph grammars provide another mechanism for integrating local
interaction rules with robot dynamics and switching logic. They have been used to
encode decentralized coordination, local information exchange, and mode
transitions in a unified framework, with applications to coverage, modular
self-reconfiguration, and deployment
\cite{mcnew2006locally,mcnew2007solving,pickem2013complete,smith2009multi}.
Their locality makes them attractive for large-scale systems, especially when
only neighboring interactions are available. The hybrid framework discussed in
Chapter~\ref{chapter:continuous}, following \cite{guo2016hybrid}, draws on this
tradition to couple local temporal tasks with interaction constraints more
systematically than is common in purely discrete coordination methods.

\subsection{Online, contingent, and human-supervised coordination}

An important recent trend is the transition from static mission synthesis to
continual and reactive coordination. Real-time reactive task allocation and
planning, online cooperation under precedence constraints, and hierarchical
coordination under continually arriving temporal tasks illustrate this shift
\cite{chen2025real,gosrich2025online,luo2025hulk}. These methods improve
responsiveness substantially, but many still rely on re-solving a global planning
object or on coordination structures that are expensive to update when new tasks
or requests arrive frequently.

This issue is especially pronounced in contingent collaboration, where robots
exchange short-horizon requests at execution time. Service requests capture the
case in which one robot requires another robot's temporary assistance, whereas
formation requests encode temporary relative-configuration requirements for a
joint subtask. Reactive temporal-logic formulations based on GR(1) and sensory
inputs are relevant in spirit \cite{kress2007s,kress2009temporal}, but they are
usually synthesized to cover a predefined class of environmental changes rather
than non-predefined inter-robot requests generated online. The contingent service
and formation setting considered in Chapter~\ref{chapter:contingent}, following
\cite{guo2015contingent}, is therefore more closely related to execution-time
request handling than to classical offline reactive synthesis. Earlier work on
formation constraints for robotic teams and prescribed-performance control also
provides relevant background \cite{loizou2004automatic,ogren2001control,
  olfati2006flocking,bechlioulis2008robust,karayiannidis2012multi,
  bechlioulis2014robust}, although these studies do not generally integrate
formation transients with high-level contingent task exchange.

The need for execution-time adaptation also connects naturally to human
supervision. Surveys on human--swarm and human--multi-robot interaction have
shown that scalable oversight requires abstractions that preserve global
situational awareness while avoiding low-level teleoperation
\cite{kolling2016human,ferreira2021survey,dahiya2023survey}. Mixed-initiative
teaming, scaled autonomy, and interactive fleet learning further study how
supervisory effort should be allocated across a robotic team under uncertainty
\cite{gombolay2017computational,swamy2020scaledautonomy,
  hoque2022fleetdagger}. However, the joint treatment of formal temporal task
semantics, online contingent collaboration, relative-motion safety, and selective
human intervention remains comparatively underdeveloped. This gap motivates the
overall perspective adopted in this monograph.
